\section{\BenchmarkName: Task and Evaluation Protocol}
\label{sec:benchmark}

\BenchmarkName\ evaluates a complete quantitative research run, from the first model interaction to a frozen strategy's sealed assessment. The benchmark specifies the comparison contract; the Arena platform executes that contract. This separation allows research procedures to vary while preserving the information and measurement rules used to judge their outputs.

\subsection{Task, Participant, and Evaluation Unit}
\label{sec:arena_task}

We define a benchmark suite as
\begin{equation}
    q = (\mathcal{D},\mathcal{U},\mathcal{W},\mathcal{P},\mathcal{E},\mathbf{B},\mathcal{F},\mathcal{V}),
    \label{eq:arena_task}
\end{equation}
where $\mathcal{D}$ identifies immutable data releases, $\mathcal{U}$ the authorized instrument population, $\mathcal{W}$ the chronological windows and warmup rules, $\mathcal{P}$ the input profile, $\mathcal{E}$ the execution environment, $\mathbf{B}$ the resource budgets, $\mathcal{F}$ the feedback policy, and $\mathcal{V}$ the financial and scoring evaluator. All components are frozen before a matched comparison. A change to this contract defines a new benchmark setting.

A participant $a=(M,H)$ combines a registered language model $M$ with a research harness $H$. The harness determines prompt construction, candidate generation, feedback interpretation, stopping, and final selection. A replicate executes this participant under $q$ with a registered seed identifier and its own budget. An episode contains the prescribed replicates; their complete records constitute the unit of comparison. Harness changes, including different prior memories or candidate-selection policies, define participant variants under the same suite.

The platform supports three input profiles: authorized raw market data, fifteen-second order-book observations, and anonymous numerical features. They expose different information and preprocessing requirements and are evaluated as distinct conditions. Strategies return cross-sectional stock scores for authorized signal dates and may declare constituent features. Optional training uses an earlier registered interval; the resulting model artifact is frozen and reused in research and final evaluation. Both participant-defined scoring programs and trusted linear combinations of declared features are supported. The latter provide directly verifiable feature-use evidence for diversity assessment.

\subsection{Research Interaction and Resource Constraints}
\label{sec:arena_interaction}

The harness alternates between model interactions and admitted research evaluations. Each completed evaluation returns the projection permitted by $\mathcal{F}$, including selected metrics, validity and coverage information, execution summaries, and remaining resources. Metric precision and available diagnostics are fixed by the suite. Private signal matrices, quantitative execution logs, fitted artifacts, and trading records are withheld from the research feedback channel.

The current budget profile distinguishes model tokens, admitted research evaluations, and execution CPU usage. Let $u_j$ be the reported token usage of model operation $j$, $n_{\mathrm{adm}}$ the number of admitted research evaluations, and $c_\ell$ the measured CPU usage of execution $\ell$. The corresponding usage vector is
\begin{equation}
    \mathbf{b} = \left(\sum_j u_j,\ n_{\mathrm{adm}},\ \sum_\ell c_\ell\right),
    \qquad \mathbf{b} \preceq \mathbf{B}
    \quad\text{for budget-eligible runs}.
    \label{eq:budgets}
\end{equation}
These sums apply when the required usage records are known; unresolved usage retains an explicit unavailable state. Admission also accounts for outstanding reservations and protects CPU capacity for final evaluation. Reservations are released or settled against observations and do not become measured consumption by definition. CPU accounting covers registered executions rather than all shared-service overhead. Token ceilings do not imply equal monetary expenditure or equal effective computation across models.

A malformed model reply consumes its observed tokens. A strategy rejected before evaluation admission does not consume an admitted evaluation; an accepted evaluation retains its evaluation charge if execution subsequently fails. Data requests and bytes remain metered under the current service-limit profile, with access and concurrency controls; historical profiles retain their own cumulative resource definitions.

Three participation modes establish different research conditions. \emph{One-shot} returns one complete strategy and freezes it, without research feedback, model-assisted repair, or additional training; platform documents may be read on request, and those reads are charged as ordinary model calls. \emph{Vanilla} uses a model--feedback conversation: the model may finish with a completed candidate, and the last completed candidate with feedback is the declared fallback when the budget ends. \emph{External harnesses} implement their own search and selection rules through the same authorized services. Their external knowledge, memory initialization, repair policy, and selection rule form part of the participant specification.

\subsection{Temporal Access and Final Commitment}
\label{sec:arena_contract}

Training, research, and final accounting windows are registered chronologically. Authorized warmup supplies historical context without extending the accounting interval. For a signal date $d$, the executable strategy receives only observations in the authorized data release, instrument and field scope, and availability interval ending at $\kappa(d)$. The cutoff $\kappa(d)$ is registered with the suite and is not disclosed in the public contract. Data permissions are enforced by the execution service rather than inferred from dates supplied by participant code.

After research, each replicate commits one strategy. Its source, runtime, fitted artifact when present, and declared feature configuration are fixed before sealed assessment. The harness controls the selection; the platform validates and freezes the commitment and runs the final evaluator. Final-window outcomes are excluded from candidate revision within that replicate. Further evaluation under another contract produces a separate record and preserves the original result.

The daily protocol uses registered static Shenzhen A-share factor columns. Missing observations are retained, and no daily index-constituent mask is imposed. Population construction, delisted securities, and data-publication history must therefore be disclosed alongside the suite. Runtime authorization controls access during the experiment; historical knowledge in model pretraining or participant priors remains a separate source of possible contamination. Repeated public leaderboard access also requires a campaign-level disclosure policy, as discussed in Appendix~\ref{app:boundary}.

\subsection{Daily Financial Evaluation}
\label{sec:arena_returns}
% DRAFT (Claude, factual skeleton): prose to be rewritten by Gemini; keep the
% wording rules: "bought on T, sold on T+1 with an algorithmic order"; slippage
% "computed from order-book limits and market impact, not a fixed value".

Index consecutive trading sessions by $t$. Scores $s_{i,t-1}$ are formed using information authorized through session $t-1$. A selected cohort is bought on trading day $t$ and sold on day $t+1$ with an algorithmic order. Slippage is computed from order-book limits and market impact; it is not a fixed value. Limit-up stocks are not bought, halted or limit-down positions are sold once they become tradable, and no replacement stocks are selected. An authorized predecessor signal can support the first entry day; an absent predecessor is not reconstructed from later data, and missing source rows produce no signal.

\paragraph{Selection.}
Let $\mathcal{I}_t$ contain the registered instruments with finite scores for entry $t$, ranked in descending order. For a valid signal row, let $s_{[k],t-1}$ be its $k$-th largest score and define
\begin{equation}
    k_t=\min(K,|\mathcal{I}_t|),\qquad
    \mathcal{S}^{K}_t=\{i\in\mathcal{I}_t:s_{i,t-1}\ge s_{[k_t],t-1}\}.
    \label{eq:topk}
\end{equation}
This definition includes every boundary tie, so the selected count may exceed $K$. Selection precedes the observation of forward prices and does not replace securities whose returns later prove unavailable. The primary basket uses $K=75$; $K\in\{200,400\}$ are retained as diagnostics.

\paragraph{Research returns.}
Let $y_{i,t}$ be the realized raw-price return of stock $i$ over the cohort's holding period and $\mathcal{A}^{K}_t\subseteq\mathcal{S}^{K}_t$ the selected stocks with a valid return. The cohort research return $r^K_t$ is the equal-weight mean of $y_{i,t}$ over $\mathcal{A}^K_t$, net of trading costs. Cohorts overlap and returns accumulate additively:
\begin{align}
    R^K &= \sum_{t\in\mathcal{J}_K}r^K_t,
    &\operatorname{ARR}^{K} &= \frac{243}{|\mathcal{J}_K|}\sum_{t\in\mathcal{J}_K}r^K_t,\label{eq:research_arr}\\
    \Delta\operatorname{ARR}^{K} &= \frac{243}{|\mathcal{J}^{\cap}_K|}
       \sum_{t\in\mathcal{J}^{\cap}_K}(r^K_t-r^{\mathrm{mkt}}_t),
       \label{eq:excess_arr}
\end{align}
where $\mathcal{J}_K$ are entry sessions with an available return, $r^{\mathrm{mkt}}_t$ is the full-market Shenzhen A-share reference under the same convention, and $\mathcal{J}^{\cap}_K$ are paired sessions. An empty observable set yields an unavailable return rather than zero. Signed additive drawdown is
\begin{equation}
    \operatorname{DD}^{K}
    =\min_{0\le j\le |\mathcal{J}_K|}
       \left(P^K_j-\max_{0\le h\le j}P^K_h\right)\le0,
    \label{eq:drawdown}
\end{equation}
with $P^K_j$ the cumulative sum through the $j$-th available cohort. The long-only Sharpe ratio is $\mathrm{SR}_L=\sqrt{243}\,\overline{r^{75}}/\operatorname{sd}(r^{75})$. The long--short Sharpe ratio $\mathrm{SR}_{LS}$ applies the same annualization to the daily raw return of the top score decile minus the bottom decile; it measures the signal's discriminative power, not a tradable return. Because adjacent cohorts overlap, both ratios are descriptive. A funded cash account is simulated and reported separately from research returns.

\subsection{Scoring}
\label{sec:scoring}
\label{sec:arena_scoring}
% DRAFT (Claude): formulas are final; prose to be rewritten by Gemini.

On each valid entry date, RankIC is the cross-sectional rank correlation between $s_{i,t-1}$ and $y_{i,t}$; $\overline{\mathrm{RankIC}}$ is its daily mean and RankICIR the mean divided by the sample standard deviation. Turnover $\mathrm{TO}\in[0,1]$ is the mean daily fraction of the Top75 basket that is replaced. With $x^+=\max(0,x)$, the in-sample (2022--2023) quality score is the weighted geometric mean
\begin{equation}
    Q=\Big(\sqrt{\overline{\mathrm{RankIC}}^{+}}\Big)^{w_1}
      \Big(\tfrac{A^+}{1+A^+}\Big)^{w_2}
      \Big(\tfrac{\mathrm{SR}_L^+}{1+\mathrm{SR}_L^+}\Big)^{w_3}
      \Big(\tfrac{\mathrm{SR}_{LS}^+}{1+\mathrm{SR}_{LS}^+}\Big)^{w_4}
      \Big(\tfrac{1}{1+\mathrm{TO}}\Big)^{w_5},
    \qquad A=\mathrm{RankICIR}\cdot\sqrt{243},
    \label{eq:quality}
\end{equation}
with $(w_1,\dots,w_5)=(0.225,0.225,0.225,0.225,0.1)$. Each term lies in $[0,1]$ with diminishing returns. Turnover enters as $1/(1+\mathrm{TO})$ so that a fully rotating basket halves this term instead of zeroing $Q$. Efficiency and diversity are
\begin{equation}
    E=Q\left(1-\frac{T}{B}\right),\qquad
    D=Q\,\frac{\min(K,32)}{32},
    \label{eq:efficiency}
\end{equation}
where $T$ is the provider-reported token usage, $B=500{,}000$ the token budget, and $K$ the number of qualified, used, and mutually independent feature clusters (Appendix~\ref{app:diversity}). The window score is
\begin{equation}
    S=\begin{cases}70Q+20E+10D & \text{one-shot and vanilla tracks},\\ 80Q+20D & \text{system and custom tracks,}\end{cases}
    \label{eq:overall}
\end{equation}
on a 0--100 scale, and the leaderboard total is
\begin{equation}
    \mathrm{Total}=0.3\,S_{\mathrm{in}}+0.7\,S_{\mathrm{out}}.
    \label{eq:total}
\end{equation}
The in-sample formula, its parameters, and all component scores are published. The out-of-sample (2024) score $S_{\mathrm{out}}$ takes into account RankIC, IC stability, long-only and long--short Sharpe ratios, turnover, factor diversity, and token usage; its combination is not disclosed and may be revised independently. Raw metrics of both windows are published.

Unavailable metrics retain reason codes and sample counts and are never replaced by zero. Failed runs remain in every denominator; averages over completed runs only are labelled as conditional statistics.
